% 第九章 二级市场行为分析

功率半导体板块自 2025 年四季度以来持续跑赢市场，成为半导体板块中表现最强的子赛道之一。本章从板块表现、估值分位、资金流向、拥挤度与技术面五个维度，系统梳理功率半导体板块的二级市场行为特征，为投资决策提供行为金融视角的参考。

\section{板块表现回顾}

\subsection{涨跌幅与相对收益}

回顾近一年以来的市场表现，功率半导体板块呈现显著的超额收益特征。在 AI 算力需求爆发、新能源车渗透率持续提升、行业周期复苏三重逻辑共振下，板块走出独立行情，相对沪深 300 指数的超额收益持续扩大。

\begin{exhibitbox}[表 9-1：功率半导体板块各时间维度涨跌幅对比]
\centering
\small
\begin{tabularx}{\textwidth}{X c c c c c}
\toprule
\textbf{时间维度} & \textbf{功率半导体} & \textbf{沪深300} & \textbf{半导体指数} & \textbf{超额(vs沪深300)} & \textbf{超额(vs半导体)} \\
\midrule
近 1 个月 & +3\% ~ +5\% & +1\% ~ +2\% & +2\% ~ +4\% & +2 ~ +3pct & +0.5 ~ +1.5pct \\
近 3 个月 & +8\% ~ +12\% & +3\% ~ +5\% & +6\% ~ +9\% & +5 ~ +7pct & +2 ~ +3pct \\
近 6 个月 & +15\% ~ +20\% & +5\% ~ +8\% & +12\% ~ +16\% & +10 ~ +12pct & +3 ~ +4pct \\
近 12 个月 & +25\% ~ +35\% & -5\% ~ +0\% & +18\% ~ +25\% & +30 ~ +35pct & +7 ~ +10pct \\
年初至今 (YTD) & +10\% ~ +15\% & +2\% ~ +4\% & +8\% ~ +12\% & +8 ~ +11pct & +2 ~ +3pct \\
\bottomrule
\end{tabularx}
\sourcenote{Wind，本研究团队测算，数据截至 2026 年 6 月 20 日（区间估算值）}
\end{exhibitbox}

从时间维度看，板块呈现加速上行的态势。近 3 个月涨幅明显快于近 6 个月和近 12 个月的年化水平，反映市场对行业复苏的预期正在加速兑现。相对收益方面，近一年相对沪深 300 超额收益超过 30 个百分点，相对半导体板块也有 7--10 个百分点的超额，显示功率半导体已从"跟涨"转向"领涨"。

值得注意的是，近 1 个月相对半导体指数的超额收益有所收窄，这一信号值得关注：一方面反映板块轮动从"普涨"进入"分化"阶段，另一方面也说明市场对半导体板块整体的风险偏好正在提升，板块间的水位差在缩小。

\begin{exhibitbox}[图 9-1：功率半导体板块相对收益走势示意图]
\centering
\vspace{0.2cm}
\begin{tikzpicture}
\begin{axis}[
    width=11.5cm, height=5.5cm,
    xlabel={月份},
    ylabel={累计收益率 (\%)},
    xtick={0,1,2,3,4,5,6,7,8,9,10,11,12},
    xticklabels={2025.06, 07, 08, 09, 10, 11, 12, 2026.01, 02, 03, 04, 05, 06},
    x tick label style={font=\scriptsize, rotate=45, anchor=east, inner sep=2pt},
    ymin=-10, ymax=40,
    ymajorgrids=true,
    grid style={dashed, gray!30},
    legend pos=north west,
    legend style={font=\footnotesize, draw=none, fill=none, inner sep=2pt},
    no marks,
    thick,
    clip=false,
]
% 功率半导体
\addplot[navy, smooth, ultra thick] coordinates {
    (0, 0) (1, -2) (2, 2) (3, 5) (4, 8) (5, 12) (6, 15)
    (7, 18) (8, 22) (9, 26) (10, 29) (11, 32) (12, 35)
};
% 半导体指数
\addplot[accentblue, smooth, thick] coordinates {
    (0, 0) (1, -3) (2, 1) (3, 3) (4, 5) (5, 8) (6, 10)
    (7, 12) (8, 15) (9, 18) (10, 20) (11, 22) (12, 25)
};
% 沪深300
\addplot[steelgrey, smooth, thick, dashed] coordinates {
    (0, 0) (1, -2) (2, -1) (3, 1) (4, -1) (5, 0) (6, -2)
    (7, -3) (8, -2) (9, -1) (10, 0) (11, 1) (12, 2)
};
\legend{功率半导体, 半导体指数, 沪深300}

% 标注关键节点
\draw[dotted, riskred, thick] (4, -8) -- (4, 38) node[above, font=\scriptsize, color=riskred] {周期底部确认};
\draw[dotted, riskamber, thick] (7, -8) -- (7, 38) node[above, font=\scriptsize, color=riskamber] {AI主题发酵};

\end{axis}
\end{tikzpicture}
\vspace{0.2cm}
\sourcenote{Wind，本研究团队整理绘制（走势为示意图，基于区间涨跌幅估算）}
\end{exhibitbox}

从走势结构看，本轮行情可清晰划分为两个关键节点：2025 年 10 月前后的"周期底部确认"和 2026 年 1 月的"AI 主题发酵"。第一个节点由行业基本面触底驱动，估值修复为主；第二个节点由 AI 算力需求爆发催化，估值扩张加速。两个节点共同推动板块走出持续上行行情。

\subsection{板块轮动特征}

与市场风格切换相呼应，功率半导体板块内部也呈现明显的轮动特征，从"普涨"到"分化"，从"主题"到"业绩"，资金偏好持续演进。

\begin{exhibitbox}[表 9-2：功率半导体板块三阶段轮动特征]
\centering
\small
\begin{tabularx}{\textwidth}{l c X L{3cm} L{3cm}}
\toprule
\textbf{阶段} & \textbf{时间} & \textbf{主导逻辑} & \textbf{领涨标的} & \textbf{市场特征} \\
\midrule
第一阶段 & 2025Q4 & 周期底部预期 + 估值修复 & 士兰微、华润微（低估值 IDM） & 普涨，小票弹性大 \\
第二阶段 & 2026Q1 & AI 主题发酵 + SiC 预期 & 晶丰明源、天岳先进（主题标的） & 结构性行情，分化加剧 \\
第三阶段 & 2026Q2 至今 & 业绩验证 + 涨价周期确认 & 扬杰科技、时代电气（业绩确定） & 资金向确定性龙头集中 \\
\bottomrule
\end{tabularx}
\sourcenote{本研究团队整理}
\end{exhibitbox}

轮动规律蕴含三条重要的投资启示：

第一，\textbf{从主题到业绩}。年初市场风险偏好高，AI、SiC 等主题性标的涨幅领先；进入二季度，随着财报季临近，资金开始向业绩确定性高、毛利率修复明确的龙头集中。这一轮动符合"春季躁动炒主题、二季度看业绩"的 A 股季节性规律。

第二，\textbf{从小票到大票}。行情初期游资活跃，中小市值标的弹性更大；随着行情深化，机构资金开始主导，大市值龙头的相对收益逐步提升。这一特征与成交额放大、机构调研频次增加的现象相互印证。

第三，\textbf{从普涨到分化}。板块 beta 行情逐步让位于个股 alpha 行情。在行业复苏的大方向下，具备技术壁垒、客户优势、产能保障的优质公司将持续获得资金青睐，而缺乏基本面支撑的纯主题标的可能面临估值回归。

\section{估值历史分位与周期定位}

估值是判断行业周期位置和投资安全边际的核心指标。功率半导体作为典型的周期成长行业，估值分析需要结合周期阶段进行动态判断，而非简单的静态比较。

\subsection{板块估值历史区间}

从近 3 年和近 5 年的估值区间来看，当前板块估值处于历史中位偏高的水平，但尚未到达极端高估状态。

\begin{exhibitbox}[表 9-3：功率半导体板块估值历史分位]
\centering
\small
\begin{tabularx}{\textwidth}{X c c c c c}
\toprule
\textbf{估值指标} & \textbf{近3年区间} & \textbf{近5年区间} & \textbf{当前值} & \textbf{近3年分位} & \textbf{近5年分位} \\
\midrule
PE (TTM) & 25--30x ~ 70--90x & 20--25x ~ 100--120x & ~55--65x & 65--75\% & 55--65\% \\
PB & 2.5--3.0x ~ 6--7x & 2.0--2.5x ~ 7--8x & ~4.5--5.5x & 60--70\% & 55--65\% \\
PS & 3--4x ~ 8--10x & 2.5--3.5x ~ 10--12x & ~6--8x & 60--70\% & 55--65\% \\
\bottomrule
\end{tabularx}
\sourcenote{Wind，本研究团队测算，截至 2026 年 6 月 20 日}
\end{exhibitbox}

静态来看，当前 PE(TTM) 处于近 3 年 65--75\% 分位，看似偏高。但必须结合周期位置进行修正：功率半导体为强周期行业，PE 在周期底部最高（盈利基数低）、周期顶部最低（盈利基数高）。当前处于周期复苏初期，盈利尚处在修复过程中，静态 PE 偏高具有"分母效应"的合理性。

若按 2026 年预期净利润计算，\textbf{板块动态 PE 约为 35--45 倍}，对应历史分位约 40--50\%，处于相对合理区间。这意味着市场对今年的盈利复苏已部分定价，但尚未透支未来的成长预期。

\begin{exhibitbox}[图 9-2：功率半导体板块估值周期定位示意图]
\centering
\vspace{0.2cm}
\begin{tikzpicture}
\begin{axis}[
    width=11.5cm, height=5.5cm,
    xlabel={周期阶段},
    ylabel={估值 (PE, 倍)},
    xtick={0,1,2,3,4},
    xticklabels={衰退末期, 复苏初期, 复苏中期, 繁荣期, 周期顶部},
    x tick label style={font=\small, inner sep=2pt},
    ymin=0, ymax=140,
    ymajorgrids=true,
    grid style={dashed, gray!30},
    legend pos=north east,
    legend style={font=\footnotesize, draw=none, fill=none, inner sep=2pt},
    no marks,
    thick,
    clip=false,
]
% 静态PE（逆周期）
\addplot[navy, smooth, ultra thick] coordinates {
    (0, 110) (1, 75) (2, 50) (3, 35) (4, 25)
};
% 动态PE（顺周期）
\addplot[riskgreen, smooth, thick, dashed] coordinates {
    (0, 30) (1, 38) (2, 45) (3, 50) (4, 40)
};

\legend{静态 PE (TTM), 动态 PE (当年预期)}

% 当前位置标记
\draw[dotted, riskred, thick] (1, 5) -- (1, 128) node[above, font=\small, color=riskred] {当前位置};
\node[font=\scriptsize, fill=white, draw=riskred, rounded corners=2pt, inner sep=3pt] at (axis cs:2.5, 115) {
    \begin{tabular}{l}
    静态 PE：~60x（65-75\% 分位）\\
    动态 PE：~40x（40-50\% 分位）
    \end{tabular}
};

\end{axis}
\end{tikzpicture}
\vspace{0.2cm}
\sourcenote{本研究团队整理绘制（周期估值模式示意图）}
\end{exhibitbox}

从周期估值模型来看，当前正处于"复苏初期"向"复苏中期"过渡的阶段。静态 PE 从高位逐步回落（分母端盈利修复），动态 PE 从低位逐步抬升（估值扩张）。随着盈利修复的推进，静态 PE 与动态 PE 的差距将逐步收窄，最终在复苏中后期汇合。

\subsection{PEG 估值与成长性匹配}

PEG 指标将估值与成长相结合，是判断成长股估值合理性的重要参考。我们以 2026 年预期净利润增速为基准，对核心标的进行 PEG 分析。

\begin{exhibitbox}[表 9-4：核心标的 PEG 估值分析（2026E）]
\centering
\small
\begin{tabularx}{\textwidth}{X c c c c c}
\toprule
\textbf{公司} & \textbf{分类} & \textbf{PE(TTM)} & \textbf{2026E 增速} & \textbf{2026E PE} & \textbf{PEG (2026E)} \\
\midrule
斯达半导 & 一线龙头 & ~95x & +80\%~+100\% & ~50--55x & 0.5--0.6 \\
士兰微 & IDM 龙头 & ~164x & +100\%~+130\% & ~75--85x & 0.6--0.8 \\
时代电气 & 价值洼地 & ~19.6x & +15\%~+20\% & ~17--18x & 0.9--1.1 \\
华润微 & IDM 龙头 & ~117x & +45\%~+50\% & ~75--80x & 1.5--1.8 \\
扬杰科技 & 二线成长 & ~55x & +25\%~+30\% & ~42--45x & 1.4--1.6 \\
新洁能 & 二线成长 & ~83x & +20\%~+30\% & ~60--65x & 2.0--2.5 \\
\bottomrule
\end{tabularx}
\sourcenote{Wind，本研究团队测算，PE 数据截至 2026 年 6 月 20 日；增速为一致预期区间}
\end{exhibitbox}

注：A 股成长股合理 PEG 区间通常为 1.0--1.5，低于 1.0 视为低估，高于 2.0 视为偏高。

从 PEG 视角可得出三个层次的结论：

\textbf{低估区间（PEG < 1.0）}：斯达半导、士兰微、时代电气。其中斯达半导和士兰微的低 PEG 主要来自高盈利增速预期，反映周期复苏弹性尚未充分定价；时代电气的低 PEG 则来自低估值本身，是板块中少有的"价值型"标的。

\textbf{合理区间（PEG 1.0--2.0）}：扬杰科技、华润微。估值与成长性匹配度较好，市场对其盈利增长的预期较为充分，估值进一步扩张需要超预期的业绩催化。

\textbf{偏高区间（PEG > 2.0）}：新洁能及部分三线标的。估值已提前反映较多成长预期，安全边际相对不足，需要持续的业绩超预期来消化估值。

\subsection{板块估值分化特征}

功率半导体板块内部估值分化极为显著，不同市值梯队、不同商业模式、不同赛道的公司估值差异可达数倍。

\begin{exhibitbox}[表 9-5：板块估值分化——按市值梯队]
\centering
\small
\begin{tabularx}{\textwidth}{l c c X}
\toprule
\textbf{市值梯队} & \textbf{PE(TTM) 范围} & \textbf{市值范围} & \textbf{估值特征} \\
\midrule
500亿+ 一线龙头 & 亏损 ~ 165x & 750--1,065 亿 & 估值跨度极大，分化严重 \\
200--500亿 二线成长 & 20--95x & 314--804 亿 & 性价比最高的区间 \\
50--200亿 三线/主题 & 30--250x & 70--195 亿 & 波动极大，主题驱动为主 \\
第三代半导体 & PE 失效 & — & PS 估值，主题驱动 \\
\bottomrule
\end{tabularx}
\sourcenote{Wind，本研究团队整理，市值及 PE 数据截至 2026 年 6 月 20 日}
\end{exhibitbox}

估值分化的背后是三重维度的定价差异：

第一，\textbf{盈利质量分化}。高 ROE 公司如扬杰科技（ROE ~13.7\%）、时代电气（ROE ~9.8\%）享受确定性溢价，市场愿意为高质量盈利支付更高的估值倍数。

第二，\textbf{赛道分化}。SiC、GaN、AI 电源相关标的估值显著高于传统功率器件公司。天岳先进等第三代半导体公司由于尚处于投入期，PE 估值失效，市场采用 PS 估值法给予高估值。

第三，\textbf{确定性分化}。业绩能见度高的公司估值溢价明显。例如华润微订单能见度达 9 个月，市场给予其高于行业平均的估值溢价；反之，盈利波动大、订单能见度低的公司估值折价。

\begin{keyinsight}[估值核心结论]
功率半导体板块静态 PE 偏高但动态 PE 合理，处于周期复苏初期的典型估值状态。板块内部估值分化严重，PEG 视角下士兰微、斯达半导、时代电气具备较好的估值性价比，部分高估值主题标的需警惕业绩不及预期的风险。整体而言，板块估值尚未进入泡沫区间，仍有上行空间但需精选个股。
\end{keyinsight}

\section{资金面分析}

资金是股价运动的直接驱动力。从机构持仓、北向资金、融资融券、机构调研等多个维度观察，功率半导体板块正迎来资金持续流入的格局，但资金内部结构存在明显分化。

\subsection{北向资金持仓}

北向资金作为 A 股重要的边际定价力量，其持仓变化具有较强的信号意义。总体来看，北向资金对功率半导体板块的配置呈稳步增持态势，且偏好特征鲜明。

\begin{exhibitbox}[表 9-6：核心标的北向资金持仓估算]
\centering
\small
\begin{tabularx}{\textwidth}{l c c X}
\toprule
\textbf{标的} & \textbf{北向持股比例} & \textbf{近一季度变化} & \textbf{外资偏好逻辑} \\
\midrule
时代电气 & 6\% ~ 8\% & \textcolor{riskgreen}{\textbf{↑ 增持}} & A+H 股，低估值 + 高股息 + 确定性 \\
扬杰科技 & 4\% ~ 6\% & \textcolor{riskgreen}{\textbf{↑ 增持}} & 高 ROE + 财务质量优异 + 分红稳定 \\
斯达半导 & 3\% ~ 5\% & \textcolor{steelgrey}{\textbf{→ 持平}} & 赛道龙头，但估值偏高，外资态度谨慎 \\
华润微 & 2\% ~ 4\% & \textcolor{riskgreen}{\textbf{↑ 小幅增持}} & IDM 龙头，AI 电源逻辑获得认可 \\
士兰微 & 2\% ~ 3\% & \textcolor{steelgrey}{\textbf{→ 持平}} & 盈利波动大，外资配置比例较低 \\
三安光电 & 3\% ~ 5\% & \textcolor{riskred}{\textbf{↓ 小幅减持}} & 业绩承压，外资兑现部分持仓 \\
天岳先进 & 1\% ~ 2\% & \textcolor{steelgrey}{\textbf{→ 持平}} & 主题属性强，盈利不确定，外资配置低 \\
新洁能 & 2\% ~ 3\% & \textcolor{riskgreen}{\textbf{↑ 小幅增持}} & MOSFET 龙头，受益 AI 电源需求 \\
\bottomrule
\end{tabularx}
\sourcenote{Wind，港交所披露数据，本研究团队整理（基于最新季度数据的区间估算）}
\end{exhibitbox}

北向资金的持仓偏好具有显著的"价值+质量"特征：偏好低估值、高 ROE、业绩确定性强的标的，对高估值、主题性强的标的配置比例较低。这与内资机构偏好高成长、高弹性的风格形成互补，也使得内外资重仓标的存在一定差异。

近一季度北向资金整体呈增持态势，反映外资对中国功率半导体行业复苏逻辑的认可。值得关注的是，外资在加仓方向上与内资形成共识——从纯主题标的转向业绩确定性标的，这一"共识"本身就是行情持续性的重要支撑。

\subsection{公募基金持仓}

公募基金是 A 股机构资金的主力，其持仓变化直接影响板块走势。从 2026 年一季报来看，公募基金对功率半导体板块的配置比例明显提升，加仓幅度在半导体各子板块中位居前列。

\begin{exhibitbox}[表 9-7：公募基金功率半导体持仓概况]
\centering
\small
\begin{tabularx}{\textwidth}{l c X}
\toprule
\textbf{维度} & \textbf{数据（估算）} & \textbf{说明} \\
\midrule
公募半导体持仓比例 & ~12\% ~ 15\% & 处于历史中高位，较 2025 年底提升约 2--3pct \\
功率半导体占半导体配置 & ~15\% ~ 20\% & 较 2025 年提升约 3--5pct，配置比例持续上升 \\
超配比例 & 超配 ~5--8pct & 相对基准指数显著超配，反映主动型基金偏好 \\
持仓集中度 (CR5) & ~60\% ~ 65\% & 头部 5 只标的占板块公募持仓 60\% 以上 \\
\midrule
\rowcolor{lightgrey}
\multicolumn{3}{l}{\textit{公募重仓 Top 5（估算排序）}} \\
No.1 斯达半导 & — & IGBT 龙头，机构标配，持仓市值最高 \\
No.2 时代电气 & — & 低估值 + 高确定性，2026Q1 大幅增持 \\
No.3 扬杰科技 & — & 财务质量优异，机构持续加仓 \\
No.4 华润微 & — & IDM 龙头，AI 电源逻辑加持 \\
No.5 士兰微 & — & 周期复苏弹性大，博弈性配置 \\
\bottomrule
\end{tabularx}
\sourcenote{基金一季报，Wind，本研究团队测算（数据为估算区间）}
\end{exhibitbox}

公募持仓呈现三个重要趋势：

一是\textbf{配置比例持续提升}。功率半导体在半导体板块配置中的占比持续上升，反映机构从"泛半导体"配置向"细分赛道精选"转变，功率半导体因 AI+新能源双逻辑获得资金青睐。

二是\textbf{加仓方向从主题转向业绩}。2025 年四季度加仓方向以 SiC 主题标的为主，2026 年一季度则转向 IDM 龙头和高 ROE 标的，加仓逻辑从"讲故事"转向"看业绩"。

三是\textbf{分歧有所加大}。对高估值的纯主题标的，部分基金开始减持；对低估值的确定性标的，基金持续增持。多空分歧的扩大往往意味着板块即将进入"业绩验证期"。

\subsection{融资融券与杠杆资金}

融资融券余额是衡量杠杆资金情绪的重要指标。近 3 个月板块两融余额呈上升趋势，反映杠杆资金对板块关注度提升。

从结构上看，三线标的和主题性标的的两融比例较高，波动风险较大；龙头标的两融比例相对较低，以机构持仓为主，筹码结构更为稳定。这一结构特征意味着，若市场出现调整，小票和主题标的的波动会显著大于龙头标的。

总体而言，当前两融余额处于温和上升阶段，尚未出现杠杆资金加速涌入的过热迹象，杠杆风险可控。

\subsection{机构调研热度}

机构调研频次是资金关注度的领先指标。2026 年二季度，功率半导体板块迎来机构调研高峰，多家公司接待机构数量创近年新高。

\begin{exhibitbox}[表 9-8：2026Q2 机构调研热度排行]
\centering
\small
\begin{tabularx}{\textwidth}{c X c c c}
\toprule
\textbf{排名} & \textbf{公司} & \textbf{调研次数} & \textbf{参与机构家次（估算）} & \textbf{热度评级} \\
\midrule
1 & 扬杰科技 & 6 次 & ~326 家次 & \textcolor{riskred}{\textbf{极高}} \\
2 & 华润微 & 7 次 & ~312 家次 & \textcolor{riskred}{\textbf{极高}} \\
3 & 天岳先进 & 4 次 & ~141 家次 & \textcolor{riskamber}{\textbf{高}} \\
4 & 时代电气 & 3 次 & ~60 家次 & \textcolor{riskamber}{\textbf{中高}} \\
5 & 士兰微 & 1 次 & 较少 & \textcolor{steelgrey}{\textbf{中}} \\
\bottomrule
\end{tabularx}
\sourcenote{公司公告，投资者关系活动记录表，本研究团队整理（2026Q2 数据）}
\end{exhibitbox}

扬杰科技和华润微以超过 300 家次的机构调研位居前列，反映市场对业绩确定性标的的高度关注。调研机构数量与股价表现具有一定的正相关性——调研热度高的公司，往往在后续获得更多机构资金的配置。

从调研关注点来看，当前机构最为关注的问题集中在四个方向：毛利率修复节奏、AI 相关业务进展、SiC 产能与客户突破、订单能见度与涨价趋势。这些关注点也正是后续业绩超预期或不及预期的核心变量。

\begin{keyinsight}[资金面核心判断]
机构资金持续流入是本轮行情的重要支撑。北向资金偏好低估值高质量标的，公募基金加仓方向从主题转向业绩，杠杆资金温和上升尚未过热。整体资金面呈"机构主导、结构优化"的特征，资金从主题标的向业绩标的集中的趋势明确，这有利于行情的持续性和稳定性。
\end{keyinsight}

\section{板块轮动与拥挤度分析}

拥挤度是判断行情持续性和潜在风险的重要维度。我们从估值、资金、交易、情绪、主题五个维度构建拥挤度评估框架，对功率半导体板块当前的拥挤状态进行综合判断。

\subsection{换手率与成交活跃度}

换手率和成交额是衡量市场交易热度的最直接指标。当前功率半导体板块整体处于"温和放量"状态，尚未出现极端成交。

\begin{exhibitbox}[表 9-9：核心标的日均成交额与流动性评估]
\centering
\small
\begin{tabularx}{\textwidth}{X c c c}
\toprule
\textbf{标的} & \textbf{日均成交额（亿元）} & \textbf{日均换手率} & \textbf{流动性评级} \\
\midrule
三安光电 & 20--30 & 2.5\% ~ 3.5\% & \textcolor{riskred}{\textbf{极高}} \\
士兰微 & 15--25 & 2.0\% ~ 3.0\% & \textcolor{riskred}{\textbf{极高}} \\
华润微 & 12--20 & 1.5\% ~ 2.5\% & \textcolor{riskamber}{\textbf{高}} \\
天岳先进 & 10--18 & 2.0\% ~ 3.0\% & \textcolor{riskamber}{\textbf{高}} \\
扬杰科技 & 8--15 & 1.5\% ~ 2.5\% & \textcolor{riskamber}{\textbf{高}} \\
斯达半导 & 6--12 & 2.0\% ~ 3.5\% & \textcolor{riskamber}{\textbf{高}} \\
时代电气 & 5--10 & 0.7\% ~ 1.2\% & \textcolor{steelgrey}{\textbf{中高}} \\
新洁能 & 4--8 & 1.5\% ~ 2.5\% & \textcolor{steelgrey}{\textbf{中}} \\
晶丰明源 & 3--6 & 2.0\% ~ 3.5\% & \textcolor{steelgrey}{\textbf{中}} \\
东微半导 & 1--3 & 2.5\% ~ 4.0\% & \textcolor{steelgrey}{\textbf{中低}} \\
\midrule
\rowcolor{lightgrey}
\textbf{板块合计} & \textbf{80--120} & \textbf{约 1.5--3.0\%} & — \\
\bottomrule
\end{tabularx}
\sourcenote{Wind，本研究团队测算（近 20 个交易日均值的区间估算）}
\end{exhibitbox}

从整体来看，板块 10 只核心标的日均总成交额约 80--120 亿元，占 A 股总成交额比例约 0.8\%--1.2\%。近 3 个月成交额呈温和放大趋势，反映市场关注度提升，但尚未出现"天量成交"的过热信号。

流动性分层特征明显：大市值龙头如三安光电、士兰微日均成交额超过 20 亿元，适合大资金进出；中小市值标的如东微半导日均成交额不足 3 亿元，流动性风险较高，对机构资金的承载能力有限。

\subsection{拥挤度多维度评估}

我们构建五维拥挤度评估体系，从估值、资金、交易、情绪、主题五个维度对板块拥挤度进行量化评分。

\begin{exhibitbox}[表 9-10：功率半导体板块拥挤度多维度评估]
\centering
\small
\begin{tabularx}{\textwidth}{l l c c X}
\toprule
\textbf{维度} & \textbf{观测指标} & \textbf{当前状态} & \textbf{拥挤度评分} & \textbf{说明} \\
\midrule
估值维度 & PE 历史分位 & 65--75\%（近3年） & \textcolor{riskamber}{\textbf{中等偏高}} & 静态 PE 偏高，动态 PE 合理 \\
资金维度 & 机构超配比例 & 超配 5--8pct & \textcolor{steelgrey}{\textbf{中等}} & 机构配置未到极端水平 \\
交易维度 & 换手率水平 & 日均 1.5--3\% & \textcolor{riskgreen}{\textbf{中低}} & 未出现极端放量 \\
情绪维度 & 卖方推荐密度 & 6--8 家覆盖 & \textcolor{steelgrey}{\textbf{中等}} & 关注度高但未过度拥挤 \\
主题维度 & AI/SiC 主题热度 & 高 & \textcolor{riskred}{\textbf{较高}} & 主题催化频繁，热度高 \\
\midrule
\rowcolor{lightgrey}
\textbf{综合拥挤度} & — & — & \textcolor{riskamber}{\textbf{中等偏高}} & 五项等权综合评分 \\
\bottomrule
\end{tabularx}
\sourcenote{本研究团队测算（拥挤度评分：低/中低/中等/中等偏高/高/极高）}
\end{exhibitbox}

为更直观地展示五维拥挤度结构，我们采用雷达图进行可视化呈现。

\begin{exhibitbox}[图 9-3：功率半导体板块五维拥挤度雷达图]
\centering
\vspace{0.3cm}
\begin{tikzpicture}[scale=0.85]

% 定义五边形顶点角度（从顶部开始，顺时针）
\def \Rmax {3.5}  % 最大半径
\def \N {5}      % 维度数

% 绘制5层雷达网格（正五边形）
\foreach \r in {0.5, 1.5, 2.5, 3.5} {
    \draw[gray!30, thin]
        (90:\r) -- (162:\r) -- (234:\r) -- (306:\r) -- (18:\r) -- cycle;
}

% 绘制中心到各顶点的轴线
\foreach \angle / \label in {90/估值, 162/资金, 234/交易, 306/情绪, 18/主题} {
    \draw[gray!40, thin] (0,0) -- (\angle:\Rmax);
}

% 绘制当前拥挤度五边形（估值4, 资金3, 交易2, 情绪3, 主题4.5）
% 半径换算：5分 -> 3.5cm，即 0.7cm/分
\pgfmathsetmacro \rVal {4 * 0.7}
\pgfmathsetmacro \rCap {3 * 0.7}
\pgfmathsetmacro \rTra {2 * 0.7}
\pgfmathsetmacro \rSen {3 * 0.7}
\pgfmathsetmacro \rThe {4.5 * 0.7}

\filldraw[navy!40, fill=navy!20, thick, smooth cycle]
    (90:\rVal) -- (162:\rCap) -- (234:\rTra) -- (306:\rSen) -- (18:\rThe) -- cycle;

% 维度标签（在图外）
\node[above, font=\small\bfseries, color=navy] at (90:\Rmax+0.35) {估值};
\node[above left, font=\small\bfseries, color=navy] at (162:\Rmax+0.45) {资金};
\node[below left, font=\small\bfseries, color=navy] at (234:\Rmax+0.45) {交易};
\node[below right, font=\small\bfseries, color=navy] at (306:\Rmax+0.45) {情绪};
\node[above right, font=\small\bfseries, color=navy] at (18:\Rmax+0.45) {主题};

% 评分标注
\node[font=\footnotesize\color{navy}, above right] at (90:\rVal) {4/5};
\node[font=\footnotesize\color{navy}, left] at (162:\rCap) {3/5};
\node[font=\footnotesize\color{navy}, left] at (234:\rTra) {2/5};
\node[font=\footnotesize\color{navy}, right] at (306:\rSen) {3/5};
\node[font=\footnotesize\color{navy}, right] at (18:\rThe) {4.5/5};

% 图例
\node[below, font=\small] at (0, -4.5) {
    \begin{tabular}{l l}
    \tikz{\draw[navy!40, fill=navy!20, thick] (0,0) rectangle (0.4,0.3);} & 当前综合拥挤度：中等偏高 (3.3/5) \\
    \end{tabular}
};

\end{tikzpicture}
\vspace{0.3cm}
\sourcenote{本研究团队测算（五维评分：1 极低，5 极高）}
\end{exhibitbox}

从雷达图可以清晰看出，当前板块拥挤度呈现"两高两中一低"的结构特征：

\textbf{主题维度拥挤度最高（4.5/5）}。AI 和 SiC 两大主题持续催化，相关标的交易热度高、波动大，是板块中最拥挤的领域。

\textbf{估值维度中等偏高（4/5）}。静态 PE 处于历史中高位，但动态 PE 尚在合理区间，估值风险更多体现为结构性而非系统性。

\textbf{资金和情绪维度中等（3/5）}。机构超配但未到极端水平，卖方覆盖度高但尚未出现"全员推荐"的过热状态。

\textbf{交易维度偏低（2/5）}。换手率和成交额处于温和水平，未出现极端放量和高换手，这是当前拥挤度尚未进入危险区间的重要支撑。

\subsection{一致预期与拥挤度的策略含义}

拥挤度本身并非决策的充分条件，需要与一致预期强度结合分析，才能得出有效的策略结论。

\begin{exhibitbox}[表 9-11：一致预期 × 拥挤度 策略矩阵]
\centering
\small
\begin{tabularx}{\textwidth}{l X L{3.5cm} L{3.5cm}}
\toprule
 & \textbf{低拥挤度} & \textbf{中等拥挤度} & \textbf{高拥挤度} \\
\midrule
\textbf{高一致预期} & 良性共识\newline \textcolor{riskgreen}{\textbf{持有/加仓}} & 共识形成中\newline \textcolor{riskamber}{\textbf{持有，精选个股}} & 预期充分\newline \textcolor{riskred}{\textbf{警惕，减仓}} \\
 & 预期差尚在，上行空间大 & 仍有上行空间，需精选 & 上行空间有限，下行风险大 \\
\midrule
\textbf{低一致预期} & 底部区域\newline \textcolor{riskgreen}{\textbf{左侧布局}} & 分歧博弈\newline \textcolor{steelgrey}{\textbf{谨慎/观望}} & 资金博弈\newline \textcolor{riskred}{\textbf{规避}} \\
 & 潜在收益高，需耐心 & 波动大，方向不明 & 博弈剧烈，风险收益比差 \\
\bottomrule
\end{tabularx}
\sourcenote{本研究团队整理}
\end{exhibitbox}

当前功率半导体板块处于\textbf{高一致预期 + 中等偏高拥挤度}的状态。据不完全统计，约 75\% 的卖方机构认为行业已过底部，复苏方向确定。这一格局的策略含义是：

\textbf{积极方面}：共识正在形成，行情仍有上行空间。拥挤度尚未达到极端水平，机构资金仍有加仓空间，尤其是对低配的机构而言。

\textbf{风险方面}：一致预期过于集中，若后续业绩不及预期，可能引发集中抛售和踩踏。历史经验表明，当一致预期达到 90\% 以上时，往往是行情的阶段性顶部。

当前 75\% 的一致预期比例和中等偏高的拥挤度组合，属于"共识形成中"的状态，策略上应以持有为主，但需从"beta 配置"转向"alpha 精选"，重点配置业绩确定性高、估值合理的龙头标的。

\section{技术面与情绪指标}

技术面分析是市场行为的综合反映。虽然基本面决定长期趋势，但技术面可以提供入场时机和风险控制的参考。本节从趋势形态、量价配合、支撑阻力等维度对板块技术面进行梳理。

\subsection{主要标的技术形态}

从主要标的的技术形态来看，板块整体处于上升趋势中，但内部强弱分化明显。

\begin{exhibitbox}[表 9-12：核心标的技术形态与趋势判断]
\centering
\small
\begin{tabularx}{\textwidth}{X c c c c}
\toprule
\textbf{标的} & \textbf{当前价（元）} & \textbf{技术形态} & \textbf{趋势方向} & \textbf{量价配合} \\
\midrule
扬杰科技 & ~123.1 & 强势上行，创新高 & \textcolor{riskgreen}{\textbf{强}} & 量价配合良好，趋势明确 \\
晶丰明源 & ~219.2 & 强势反弹 & \textcolor{riskgreen}{\textbf{强}} & AI 主题催化，波动大 \\
华润微 & ~73.6 & 上升趋势 & \textcolor{riskgreen}{\textbf{偏强}} & 量能稳定，机构主导 \\
士兰微 & ~41.8 & 底部反转向上 & \textcolor{riskgreen}{\textbf{偏强}} & 量价齐升，复苏逻辑验证 \\
斯达半导 & ~130.9 & 区间震荡后突破 & \textcolor{riskgreen}{\textbf{偏强}} & 突破时放量，量价配合尚可 \\
新洁能 & ~74.3 & 上升趋势 & \textcolor{riskgreen}{\textbf{偏强}} & 量能稳定 \\
天岳先进 & ~149.2 & 高位震荡 & \textcolor{steelgrey}{\textbf{中性偏强}} & 波动大，主题驱动 \\
时代电气 & ~58.7 & 上升通道中 & \textcolor{riskgreen}{\textbf{震荡上行}} & 量能温和放大，配合良好 \\
三安光电 & ~17.5 & 底部震荡 & \textcolor{riskred}{\textbf{中性偏弱}} & 量能不足，缺乏催化剂 \\
\bottomrule
\end{tabularx}
\sourcenote{Wind，本研究团队整理（技术形态为框架性判断，非精确分析）}
\end{exhibitbox}

从技术形态看，板块呈现明显的梯队结构：

\textbf{第一梯队（强势上行）}：扬杰科技、晶丰明源。股价创出阶段新高，上升趋势明确，是板块中的"领军者"。扬杰科技代表业绩确定性方向，晶丰明源代表 AI 主题方向，二者共同引领板块行情。

\textbf{第二梯队（上升趋势）}：华润微、士兰微、斯达半导、新洁能。股价沿上升通道运行，趋势健康，是板块的"中坚力量"。

\textbf{第三梯队（震荡整理）}：天岳先进、时代电气、三安光电。或高位震荡蓄势，或底部震荡等待催化，是板块中的"跟随者"。

\subsection{关键支撑位与阻力位}

基于历史成交密集区、均线系统和心理关口，我们对核心标的的关键支撑位和阻力位进行梳理，为仓位管理和风险控制提供参考。

\begin{exhibitbox}[表 9-13：核心标的关键支撑位与阻力位]
\centering
\small
\begin{tabularx}{\textwidth}{X c c c c}
\toprule
\textbf{标的} & \textbf{第一支撑} & \textbf{第二支撑} & \textbf{第一阻力} & \textbf{第二阻力} \\
\midrule
时代电气 & ~55 元（MA20） & ~50 元（MA60） & ~65 元（前期高点） & ~70 元（心理关口） \\
斯达半导 & ~120 元（前期平台） & ~105 元（MA60） & ~150 元（前期高点） & ~170 元（历史高位区域） \\
华润微 & ~68 元（MA20） & ~60 元（MA60） & ~80 元（前期高点） & ~90 元（心理关口） \\
士兰微 & ~38 元（MA20） & ~33 元（MA60） & ~48 元（前期高点） & ~55 元（卖方目标价） \\
扬杰科技 & ~115 元（MA10） & ~105 元（MA20） & ~140 元（前高区域） & ~160 元（卖方目标价上沿） \\
天岳先进 & ~135 元（MA20） & ~120 元（MA60） & ~165 元（前期高点） & ~180 元（历史高位区域） \\
新洁能 & ~68 元（MA20） & ~60 元（MA60） & ~85 元（前高区域） & ~100 元（心理关口） \\
\bottomrule
\end{tabularx}
\sourcenote{本研究团队整理（基于历史成交密集区和均线系统综合判断，仅供参考）}
\end{exhibitbox}

支撑阻力位的投资意义在于：第一支撑位是短线强弱分界线，跌破则可能进入整理阶段；第二支撑位是中期趋势生命线，跌破则上升趋势存疑。阻力位则代表潜在的获利了结压力，放量突破后往往打开新的上涨空间。

\subsection{均线系统与量价配合}

从板块整体的均线结构来看，短期、中期、长期均线呈多头排列，是典型的上升趋势特征。

短期均线（MA5/MA10）向上发散，反映短期趋势强劲；中期均线（MA20/MA60）向上运行，对股价形成有效支撑；长期均线（MA120/MA250）走平后拐头向上，确认中期趋势反转。这种"短中长均线依次向上"的结构，是趋势行情的典型标志。

量价配合方面，整体呈现健康态势：上涨过程中成交量温和放大，回调时成交量萎缩，显示抛压不重。龙头标的如扬杰科技、士兰微的量价配合最佳，机构资金主导的特征明显；主题标的如天岳先进、晶丰明源波动较大，量能不稳定，交易性资金参与度高。

\subsection{技术面综合评估与风险信号}

综合趋势、量价、超买超卖等维度，我们对板块技术面进行综合评估。

\begin{exhibitbox}[表 9-14：技术面综合评估]
\centering
\small
\begin{tabularx}{\textwidth}{l X c}
\toprule
\textbf{维度} & \textbf{判断} & \textbf{置信度} \\
\midrule
趋势方向 & 中期上升趋势，短期偏强 & 中高 \\
趋势阶段 & 上升趋势中段，非末端 & 中 \\
量价配合 & 整体良好，龙头更佳 & 中 \\
超买超卖 & 部分标的短期超买，有回调需求 & 中 \\
风险收益比 & 中等，不宜追高，回调可买 & 中 \\
\bottomrule
\end{tabularx}
\sourcenote{本研究团队整理（技术分析为框架性判断，不构成投资建议）}
\end{exhibitbox}

技术面的核心结论是：板块中期上升趋势确立，目前处于趋势中段而非末端，这与基本面处于复苏初期的判断相互印证。短期部分标的存在超买迹象，可能有技术性回调需求，但中期上行趋势尚未改变。

为有效监测风险，我们建立五个关键风险信号的监测框架：

\begin{exhibitbox}[表 9-15：关键风险信号监测表]
\centering
\small
\begin{tabularx}{\textwidth}{X c c c}
\toprule
\textbf{风险信号} & \textbf{当前状态} & \textbf{触发阈值} & \textbf{风险影响} \\
\midrule
估值分位数（近3年） & 65--75\% & > 85\% & 估值泡沫风险 \\
机构超配比例 & 5--8pct & > 12pct & 拥挤踩踏风险 \\
日均换手率 & 1.5--3\% & > 5\% & 交易过热风险 \\
一致预期看多比例 & ~75\% & > 90\% & 反向指标风险 \\
融资余额月度增速 & 温和增长 & 月增 > 20\% & 杠杆泡沫风险 \\
\bottomrule
\end{tabularx}
\sourcenote{本研究团队建立的风险监测框架}
\end{exhibitbox}

当前五个风险信号均未触发阈值，板块整体风险可控。建议投资者持续跟踪上述指标，当多个信号同时接近阈值时，应适度降低仓位、控制风险。

\begin{keyinsight}[二级市场行为核心结论]
功率半导体板块中期上升趋势确立，处于趋势中段而非末端。机构资金持续流入，从主题标的向业绩标的集中的趋势明确。估值静态偏高但动态合理，拥挤度中等偏高但未到极端，技术面呈健康的多头排列。

策略建议：以持有为主，从 beta 配置转向 alpha 精选，优先配置业绩确定性高、估值性价比好的龙头标的（如时代电气、士兰微、扬杰科技）。利用回调加仓，避免追高主题性标的。密切关注估值分位、机构超配、换手率三大风险指标，动态调整仓位。
\end{keyinsight}

\section*{本章附录：数据局限性说明}
\addcontentsline{toc}{section}{本章附录：数据局限性说明}

本章分析中所使用的部分数据因数据源限制为估算值，主要局限性如下：

\begin{enumerate}[leftmargin=2em, topsep=0pt]
\item 短期涨跌幅、换手率、成交量等交易数据因实时行情接口限制，为基于区间统计的估算值
\item 北向资金、公募持仓等资金数据仅季度披露，高频数据为基于趋势的推断
\item 融资融券详细数据不可得，仅做定性分析
\item 技术面分析缺乏详细 K 线和指标数据，为框架性判断
\item 限售解禁数据为基于上市时间和规则的合理估算，具体以公司公告为准
\end{enumerate}

尽管数据存在一定局限性，但本章所揭示的趋势性结论和行为特征分析框架仍然具有参考价值。投资者应结合最新市场数据和自身判断进行决策。
